\chapter*{Abstract}
\label{chap:abstract}
%
Industrial service operations generate large volumes of unstructured data in the form of
service tickets—multi-lingual email threads, attachments, and PDF documents—that are
difficult to process at scale using conventional methods. This thesis investigates whether
Large Language Models~(LLMs) can reliably extract actionable insights from such tickets
in an industrial context.

A preprocessing pipeline converts raw ticket files~(MSG, PDF, DOCX) into structured
Markdown representations. A benchmark dataset of 250 manually annotated tickets is
constructed from three ticket categories and used to evaluate ten LLM configurations,
including both open-source models served via Ollama~(GPT-OSS~20B, Gemma~3, Gemma~4,
Llama~3.1, Llama~3.2, DeepSeek-R1, Qwen~3) and proprietary models~(GPT-4o,
GPT-3.5-Turbo). Two complementary
evaluation protocols are applied: comparison against human-annotated ground-truth tags
using F1-score, and an LLM-as-judge framework that scores tag usefulness and reasoning
quality.

The results show that, with well-designed prompts, LLMs can extract structured
information~(problem, solution, item identifiers) and generate meaningful categorical
tags from multi-lingual service tickets. Rule-based regex extraction is sufficient for
recurring structured fields such as document categories, article codes, and machine
serial numbers, while LLMs are required for the narrative, context-dependent parts of
tickets. Among the evaluated models, \textsc{gpt-oss}~(20B) achieved the strongest
overall balance of tag usefulness and reasoning coherence on the benchmark set.

\bigskip
\noindent\textbf{Keywords:} Large Language Models, Service Ticket Analysis, Information
Extraction, Multi-label Classification, Ollama, Zero-shot Prompting, Industrial NLP.